\documentclass[10pt,a4paper]{amsart}
\usepackage{amsmath,amssymb,amsthm,graphicx}
\usepackage{verbatim}
\newenvironment{reference}{\comment}{\endcomment}
%\newenvironment{reference}{}{}
\newenvironment{slogan}{\comment}{\endcomment}
\newenvironment{history}{\comment}{\endcomment}

% For commutative diagrams we use Xy-pic
\usepackage[all]{xy}

% We use 2cell for 2-commutative diagrams.
\xyoption{2cell}
% Theorem environments.
%
\theoremstyle{plain}
\newtheorem{theorem}[subsection]{Theorem}
\newtheorem{proposition}[subsection]{Proposition}
\newtheorem{lemma}[subsection]{Lemma}

\theoremstyle{definition}
\newtheorem{definition}[subsection]{Definition}
\newtheorem{example}[subsection]{Example}
\newtheorem{exercise}[subsection]{Exercise}
\newtheorem{situation}[subsection]{Situation}

\theoremstyle{remark}
\newtheorem{remark}[subsection]{Remark}
\newtheorem{remarks}[subsection]{Remarks}

\numberwithin{equation}{subsection}

% Macros
%
\def\lim{\mathop{\mathrm{lim}}\nolimits}
\def\colim{\mathop{\mathrm{colim}}\nolimits}
\def\Spec{\mathop{\mathrm{Spec}}}
\def\Hom{\mathop{\mathrm{Hom}}\nolimits}
\def\Ext{\mathop{\mathrm{Ext}}\nolimits}
\def\SheafHom{\mathop{\mathcal{H}\!\mathit{om}}\nolimits}
\def\SheafExt{\mathop{\mathcal{E}\!\mathit{xt}}\nolimits}
\def\Sch{\mathit{Sch}}
\def\Mor{\mathop{\mathrm{Mor}}\nolimits}
\def\Ob{\mathop{\mathrm{Ob}}\nolimits}
\def\Sh{\mathop{\mathit{Sh}}\nolimits}
\def\NL{\mathop{N\!L}\nolimits}
\def\CH{\mathop{\mathrm{CH}}\nolimits}
\def\proetale{{pro\text{-}\acute{e}tale}}
\def\etale{{\acute{e}tale}}
\def\QCoh{\mathit{QCoh}}
\def\Ker{\mathop{\mathrm{Ker}}}
\def\Im{\mathop{\mathrm{Im}}}
\def\Coker{\mathop{\mathrm{Coker}}}
\def\Coim{\mathop{\mathrm{Coim}}}

% Boxtimes
%
\DeclareMathSymbol{\boxtimes}{\mathbin}{AMSa}{"02}

%
% Macros for moduli stacks/spaces
%
\def\QCohstack{\mathcal{QC}\!\mathit{oh}}
\def\Cohstack{\mathcal{C}\!\mathit{oh}}
\def\Spacesstack{\mathcal{S}\!\mathit{paces}}
\def\Quotfunctor{\mathrm{Quot}}
\def\Hilbfunctor{\mathrm{Hilb}}
\def\Curvesstack{\mathcal{C}\!\mathit{urves}}
\def\Polarizedstack{\mathcal{P}\!\mathit{olarized}}
\def\Complexesstack{\mathcal{C}\!\mathit{omplexes}}
% \Pic is the operator that assigns to X its picard group, usage \Pic(X)
% \Picardstack_{X/B} denotes the Picard stack of X over B
% \Picardfunctor_{X/B} denotes the Picard functor of X over B
\def\Pic{\mathop{\mathrm{Pic}}\nolimits}
\def\Picardstack{\mathcal{P}\!\mathit{ic}}
\def\Picardfunctor{\mathrm{Pic}}
\def\Deformationcategory{\mathcal{D}\!\mathit{ef}}
\usepackage{cleveref}
\crefname{originallemma}{original lemma}{original lemmas}
\Crefname{originallemma}{Original lemma}{Original lemmas}
\makeatletter
\newlabel{lemma-L-presheaf}{{00ZH}{source}{}{sourcepointer-1}{}}
\newlabel{lemma-L-presheaf@cref}{{[originallemma][0][]00ZH}{[1][0][]source}{}{}{}}
\makeatother
\title{Covering sieves are exactly the sieves whose sheafification is an isomorphism}
\author{The Stacks Project Authors}
\date{stacks-project:sites.tex:local-snapshot-20261012; GNU FDL 1.2 or later}
\begin{document}
\maketitle
\noindent\textit{Editorial note.} All delivered mathematical text is the author's own, copied from the pinned Stacks Project chapter file sites-cohomology.tex; no statement or proof step was rewritten, reordered or removed.\par\medskip
\noindent\textit{Editorial note.} The proof is the author's complete argument; every step the author writes out is inside the counted span and no closing sentence is a placeholder.\par\medskip
\noindent\textit{Editorial note.} Every labelled reference occurring inside the counted spans is carried in assembly.external\_references as a non-load-bearing navigation pointer (same-chapter ones at their definition site, cross-chapter ones at the invocation inside the counted span), so no load-bearing mathematics is replaced by a stub. No bibliography citation occurs inside the counted spans.\par\medskip
\noindent\textit{Editorial note.} The checked chapter list for same-name collisions is the set of stacks chapters staged in this repository: chow, cohomology, divisors, etale-cohomology, flat, more-algebra, more-morphisms, morphisms, spaces-more-morphisms, varieties together with this card's chapter sites-cohomology.\par\medskip
\noindent\textit{Editorial note.} The preamble is the chapter's own retained minimal preamble (sources/stacks-sites-cohomology/retained-preamble.tex), a verbatim excerpt of the native preamble: the Xy-pic setup and the native theorem declarations with the mathematical macros.\par\medskip
\noindent\textit{Allowed standard prerequisites (editorial).}\begin{itemize}
\item Sheafification as two applications of the plus construction. Same chapter, Theorem theorem-L-topology and Lemma lemma-L-presheaf. Identifies the sheafification with two plus constructions and gives injectivity of the map between sheafifications.
\item The axioms of a topology. Same chapter, Definition definition-topology. Supplies the transitivity axiom used in the converse direction.
\end{itemize}\medskip
\ifdefined\hypertarget\hypertarget{sourcepointer-1}{}\fi
\noindent\textit{Source reference.} Original lemma 00ZH: sources/stacks-sites/pinned-original-source.tex, source lines 11917--11917. Navigation pointer invoked at proof line 11917; the counted span contains the pointer only and the cited result is not reproduced inside it, so the pointer is registered as a non-load-bearing navigation stub in the established form for Stacks material. No bibliography citation is involved.\par\medskip
\begin{lemma}
\label{lemma-sieve-sheafification}
Let $\mathcal{C}$ be a category endowed with a topology $J$.
Let $U$ be an object of $\mathcal{C}$.
Let $S$ be a sieve on $U$. The following are equivalent
\begin{enumerate}
\item The sieve $S$ is a covering sieve.
\item The sheafification $S^\# \to h_U^\#$
of the map $S \to h_U$ is an isomorphism.
\end{enumerate}
\end{lemma}

\begin{proof}
First we make a couple of general remarks.
We will use that $S^\# = LLS$, and $h_U^\# = LLh_U$.
In particular, by Lemma \ref{lemma-L-presheaf}, we see that
$S^\# \to h_U^\#$ is injective. Note that
$\text{id}_U \in h_U(U)$. Hence it gives rise to
sections of $Lh_U$ and $h_U^\# = LLh_U$ over $U$ which
we will also denote $\text{id}_U$.

\medskip\noindent
Suppose $S$ is a covering sieve. It clearly suffices to
find a morphism $h_U \to S^\#$ such that the composition
$h_U \to h_U^\#$ is the canonical map. To find such a map
it suffices to find a section $s \in S^\#(U)$ which restricts
to $\text{id}_U$. But since $S$ is a
covering sieve, the element
$\text{id}_S \in \Mor_{\textit{PSh}(\mathcal{C})}(S, S)$
gives rise to a section of $LS$ over $U$ which restricts to
$\text{id}_U$ in $Lh_U$. Hence we win.

\medskip\noindent
Suppose that $S^\# \to h_U^\#$ is an isomorphism.
Let $1 \in S^\#(U)$ be the element corresponding to
$\text{id}_U$ in $h_U^\#(U)$. Because $S^\# = LLS$
there exists a covering sieve $S'$ on $U$ such that
$1$ comes from a
$$
\varphi \in \Mor_{\textit{PSh}(\mathcal{C})}(S', LS).
$$
This in turn means that for every $\alpha : V \to U$,
$\alpha\in S'(V)$ there exists a covering sieve $S_{V, \alpha}$
on $V$ such that $\varphi(\alpha)$ corresponds to
a morphism of presheaves $S_{V, \alpha} \to S$. In other words
$S_{V, \alpha}$ is contained in $S \times_U V$. By the second
axiom of a topology we see that $S$ is a covering sieve.
\end{proof}

\end{document}
